% paper/sections/07e_viscous_3layer.tex
%
% §7.6  粘性項：3 層アーキテクチャと μ∇²u 禁止（SP-K §3–§5）
%        変粘性二相流で μ∇²u が物理的に誤りである理由と,
%        Layer A/B/C 分離による正しい離散化.
%
% ═══════════════════════════════════════════════════════════════════════════════
\subsection{粘性項：3 層アーキテクチャと変粘性対応}
\label{sec:viscous_3layer}
% ═══════════════════════════════════════════════════════════════════════════════

SP-K \cite{SPK} は変粘性二相流の粘性項を
\textbf{Layer A（速度勾配生成）・Layer B（応力テンソル組立）・
Layer C（保存形応力ダイバージェンス）}の 3 層に分解し，
各層の独立性を維持することで
$[\tau_{xy}]_\Gamma=0$ のせん断応力連続条件を自然に保証する離散化を示す．

% -------------------------------------------------------------------------------
\subsubsection{なぜ $\mu\nabla^2 u$ が誤りか}
\label{sec:viscous_mu_lap_forbidden}
% -------------------------------------------------------------------------------

非圧縮条件 $\bnabla\cdot\mathbf{u}=0$ 下で：
\begin{equation}
  \bnabla\cdot(2\mu\mathbf{D})
  = \mu\nabla^2\mathbf{u} + 2\mathbf{D}\cdot\bnabla\mu,
  \label{eq:viscous_exact_decomposition}
\end{equation}
ここで $\mathbf{D} = \tfrac{1}{2}(\bnabla\mathbf{u} + \bnabla\mathbf{u}^\top)$
は歪速度テンソル．
第 2 項 $2\mathbf{D}\cdot\bnabla\mu$ は\textbf{交差微分寄与}で，
$\bnabla\mu\neq 0$ のとき常に非零．
二相流では $\mu$ が界面で跳躍（水-空気で $\mu_l/\mu_g\approx 55$）するため
$\bnabla\mu$ は $\delta$ 関数を含み，交差微分寄与は\textbf{特異で無視不可}．

\noindent\textbf{$\mu\nabla^2\mathbf{u}$ を使うと}：
\begin{itemize}
  \item $\mu$ の空間変動から生じる接線力を欠落
  \item せん断応力連続 $[\tau_{xy}]_\Gamma = 0$ を満たさない
  \item 界面に沿ったスプリアス接線速度を生成
  \item 高粘性比流でエネルギー散逸が崩壊（$\mathbf{u}\cdot\mathbf{F}_\mu$ が
        正の値を取りうる）
\end{itemize}

% -------------------------------------------------------------------------------
\subsubsection{3 層アーキテクチャ}
\label{sec:viscous_three_layers}
% -------------------------------------------------------------------------------

\noindent\textbf{Layer A — 速度勾配生成}：
staggered MAC 配置（\S\ref{sec:collocate_layout}）で
4 種の勾配を各々の自然位置で評価：
\begin{align}
  (\partial_x u)_{i,j} &\quad\text{cell center（$u$ x-face から）}, \label{eq:viscous_layerA_1}\\
  (\partial_y v)_{i,j} &\quad\text{cell center（$v$ y-face から）}, \label{eq:viscous_layerA_2}\\
  (\partial_y u)_{i+1/2,j+1/2} &\quad\text{corner（$u$ x-face から）}, \label{eq:viscous_layerA_3}\\
  (\partial_x v)_{i+1/2,j+1/2} &\quad\text{corner（$v$ y-face から）}. \label{eq:viscous_layerA_4}
\end{align}
Layer A は\textbf{bulk セルのみで CCD}を使用（\S\ref{sec:viscous_ccd_placement}）．
界面帯（$|\phi|\le 3h$）では 2 次中心差分に落とす．

\noindent\textbf{Layer B — 応力テンソル組立}：
\begin{align}
  \tau_{xx,i,j} &= 2\mu_{i,j}\,(\partial_x u)_{i,j}, \quad
  \tau_{yy,i,j} = 2\mu_{i,j}\,(\partial_y v)_{i,j}, \label{eq:viscous_layerB_diag}\\
  \tau_{xy,i+1/2,j+1/2} &= \mu_{i+1/2,j+1/2}\bigl[(\partial_y u) + (\partial_x v)\bigr]_{i+1/2,j+1/2}. \label{eq:viscous_layerB_off}
\end{align}
対角成分 $\tau_{xx},\tau_{yy}$ はセル中心，非対角 $\tau_{xy}=\tau_{yx}$ は\textbf{コーナー}．
$\mu_{i+1/2,j+1/2}$ はコーナーでの 4 点平均（算術または調和）．

\noindent\textbf{Layer C — 保存形応力ダイバージェンス}：
\begin{align}
  (F_\mu)_x\big|_{i+1/2,j}
  &= \frac{\tau_{xx,i+1,j}-\tau_{xx,i,j}}{\Delta x_{i+1/2}}
  + \frac{\tau_{xy,i+1/2,j+1/2}-\tau_{xy,i+1/2,j-1/2}}{\Delta y_j}, \label{eq:viscous_layerC_x}\\
  (F_\mu)_y\big|_{i,j+1/2}
  &= \frac{\tau_{yx,i+1/2,j+1/2}-\tau_{yx,i-1/2,j+1/2}}{\Delta x_i}
  + \frac{\tau_{yy,i,j+1}-\tau_{yy,i,j}}{\Delta y_{j+1/2}}. \label{eq:viscous_layerC_y}
\end{align}
Layer B/C は常に保存形；CCD を用いるのは Layer A の bulk 部分のみ．

% -------------------------------------------------------------------------------
\subsubsection{コーナー粘性：$\tau_{xy}$ を第一級要素として扱う}
\label{sec:viscous_corner}
% -------------------------------------------------------------------------------

$\mu_{i+1/2,j+1/2}$ はコーナーで\textbf{定義}する必要がある（セル中心の内挿ではなく）：
\begin{equation}
  \mu_{i+1/2,j+1/2}^\text{arith}
  = \tfrac{1}{4}(\mu_{i,j} + \mu_{i+1,j} + \mu_{i,j+1} + \mu_{i+1,j+1}),
  \label{eq:viscous_mu_corner_arith}
\end{equation}
高粘性比（$\mu_l/\mu_g\gg 1$）では調和平均が安全：
\begin{equation}
  \mu_{i+1/2,j+1/2}^\text{harm}
  = \biggl[\tfrac{1}{4}\bigl(\mu_{i,j}^{-1} + \mu_{i+1,j}^{-1} + \mu_{i,j+1}^{-1} + \mu_{i+1,j+1}^{-1}\bigr)\biggr]^{-1}.
  \label{eq:viscous_mu_corner_harm}
\end{equation}

\noindent\textbf{共有配列原則}：
$\tau_{xy}$ は\textbf{1 回だけ計算}し，x-/y-運動量方程式の両方で同じ配列を参照する．
同じ式で別々に計算すると数値非対称が入り，離散エネルギー恒等式が崩れる．

\noindent\textbf{エネルギー恒等式}：
保存形と共有 $\tau_{xy}$ により
\begin{equation}
  \langle\mathbf{u},\mathbf{F}_\mu\rangle_h
  = -2\sum_{i,j}\mu_{i,j}\bigl[(\partial_x u)^2+(\partial_y v)^2\bigr]
  - 2\sum_{i,j}\mu_{i+1/2,j+1/2}\bigl[(\partial_y u) + (\partial_x v)\bigr]^2_{i+1/2,j+1/2}
  \le 0.
  \label{eq:viscous_energy_identity}
\end{equation}
等号は $\mathbf{u}\equiv\text{const}$ のみで成立．
共有 $\tau_{xy}$ 違反・非保存 Layer C は
この量を正にし得る（粘性項が運動エネルギーの\textbf{源}になる）．

% -------------------------------------------------------------------------------
\subsubsection{CCD 配置：bulk Layer A のみ}
\label{sec:viscous_ccd_placement}
% -------------------------------------------------------------------------------

界面帯では $u\in C^0$（速度連続）だが $\bnabla u\notin C^0$
（粘性跳躍による勾配キンク）．
CCD の $\mathcal{O}(h^6)$ 精度はステンシル支持で場が滑らかであることを前提とするため
キンクを跨ぐ CCD ステンシルは $\partial_n u$ の不連続を増幅し
$\mathcal{O}(1)$ 勾配誤差を生む．

\noindent\textbf{界面帯の定義（Definition A）}：
$|\phi|\le 3h$ を満たすセル / コーナーは界面帯．
\textbf{Definition B}：ステンシル内に異相のセルを含む場合（phase-indicator mismatch）．

\noindent\textbf{演算子割当}：
\begin{center}
\begin{tabular}{@{}lll@{}}
\hline
領域 & Layer A & Layer B / C \\
\hline
bulk（界面から遠い）& CCD & 保存形 \\
界面帯 & 2 次中心 or 片側 & 保存形 \\
\hline
\end{tabular}
\end{center}

\noindent\textbf{Phi/Psi-gated bulk Laplacian ショートカット}（SP-K §8）：
bulk では $\bnabla\mu=0$ のため \eqref{eq:viscous_exact_decomposition} は
$\mu\nabla^2\mathbf{u}$ に簡約する．
$\psi$ による領域分割で
\begin{equation*}
  L_\mu\mathbf{u}
  = \mathbb{1}_{\Omega_\text{bulk}}\cdot\mu\Delta_h^{\CCD}\mathbf{u}
  + \mathbb{1}_{\Omega_\Gamma}\cdot\bnabla\cdot(2\mu\mathbf{D}_h)
\end{equation*}
で bulk パスを軽量化できる．
$\Omega_\Gamma = \{\psi\in(\varepsilon_\psi, 1-\varepsilon_\psi)\}$
で $\varepsilon_\psi\approx 0.01$．

% -------------------------------------------------------------------------------
\subsubsection{Crank--Nicolson 粘性 Helmholtz とデフェクト補正}
\label{sec:viscous_cn_defect}
% -------------------------------------------------------------------------------

Crank--Nicolson の粘性 Helmholtz 系：
\begin{equation}
  \Bigl(I - \tfrac{\Delta t}{2\rho}L_H\Bigr)\mathbf{u}^*
  = \text{RHS},
  \qquad L_H = \bnabla\cdot(2\mu\mathbf{D}_h^{(H)}),
  \label{eq:viscous_cn_helmholtz}
\end{equation}
ここで $L_H$ は CCD bulk + 2 次界面帯の混合演算子．
コンパクト LHS は前処理が困難なため Krylov 直接解は高コスト．

\noindent\textbf{デフェクト補正分裂}：
$A_H = I - (\Delta t/2\rho)L_H$，$A_L = I - (\Delta t/2\rho)L_L$
（$L_L$ は全域 2 次）として，
\begin{enumerate}
  \item $r^{(m)} = \text{RHS} - A_H\mathbf{u}^{(m)}$（高次残差）
  \item $A_L\delta\mathbf{u} = r^{(m)}$（低次内部解）
  \item $\mathbf{u}^{(m+1)} = \mathbf{u}^{(m)} + \delta\mathbf{u}$
\end{enumerate}
外側ループは FGMRES（内部解が可変），
内部 $A_L$ は PCG+AMG / ADI / LU．
$L_L$ は $L_H$ と同じ $\mu$・同じ BC・同じ帯スイッチを共有する必要があり，
素朴な $\mu_0\nabla^2$ では収束しない．

\noindent\textbf{アンチパターン}：
\begin{itemize}
  \item $\mu\nabla^2\mathbf{u}$ 形式（\S\ref{sec:viscous_mu_lap_forbidden}）
  \item Layer C で CCD 使用（保存性崩壊）
  \item $\tau_{xy}$ を x-/y- 別々に計算（対称性崩壊）
  \item コーナー $\mu$ をセル中心の 2 点内挿で代用（ステンシル非整合）
  \item 界面帯で CCD 強行（$\mathcal{O}(1)$ 勾配誤差）
\end{itemize}

% ═══════════════════════════════════════════════════════════════════════════════
